% File names are typeset with \nolinkurl (hyperref) so that they can break after '/', '.' or '_'.
% Sources: README.md (commands, CPU budget), tests/test_vrlab.py (as of Nov. 11), vrlab/algorithms.py (time.process_time),
% experiments/exp_A4_hard_instance.py (VRLAB_MAX_TIME, default 600 s; 45 s per grid run of the naive scheme),
% results/log_A4.txt and log_A4b.txt (4000 s, then 8000 s for the last A4b runs), results/A4.pkl and A4b.pkl (statuses:
% 264 stop_D and 30 diverged; 128 stop_D), results/B3_summary.json (grad_star_norm, final_D), vrlab/solvers.py,
% experiments/run_all.py (list SCRIPTS, then make_digest.py unless --quick), report/make_restated.py, report/build.py.
\paragraph{Code and data.} The code, the small cached results, the summaries and the figures are available at
\begin{center}\small
\url{https://github.com/louisDGL/Accelerating-Variance-Reduction---Optimal-Complexity-of-Finite-Sum-Optimization}
\end{center}
The repository contains the Python package \nolinkurl{vrlab} (problems, high-precision optimum, algorithms, exact
instrumentation, bounds, gap attribution), a test suite (\nolinkurl{tests/test_vrlab.py}), the experiment scripts in
\nolinkurl{experiments/}, the theory notes in \nolinkurl{theory/}, the bibliography in \nolinkurl{bibliography/}, this
report in \nolinkurl{report/}, the summaries (\nolinkurl{results/*_summary.json} and \nolinkurl{results/digest.md}) and
their analysis (\nolinkurl{results/findings.md}), the small cached runs (\nolinkurl{A3.pkl}, \nolinkurl{A4.pkl},
\nolinkurl{A4b.pkl}, \nolinkurl{B2.pkl}, \nolinkurl{GammaDelta.pkl} and \nolinkurl{validation_book.pkl} in
\nolinkurl{results/}) and the figures in \nolinkurl{figures/}. The large caches
\nolinkurl{results/A2.pkl}, \nolinkurl{results/B1.pkl} and \nolinkurl{results/B3.pkl} are not versioned and are
rebuilt by the scripts that produce them. Dependencies: numpy, scipy, matplotlib (see \nolinkurl{requirements.txt});
Python~3.14 was used.

\paragraph{Commands.} The commands of \nolinkurl{README.md}, each run from the root of the repository, install the
dependencies, run the tests and the experiments, rebuild the digest and compile the documents:
\begin{small}
\begin{verbatim}
pip install -r requirements.txt
python tests/test_vrlab.py              # or: python -m pytest tests -q
python experiments/run_all.py --quick   # smoke test (a few minutes)
python experiments/run_all.py           # full experiments (many CPU hours)
python experiments/retry_timeouts.py A4 A4b   # drop the entries that hit the CPU budget
export OPENBLAS_NUM_THREADS=1 OMP_NUM_THREADS=1 VRLAB_MAX_TIME=4000
python experiments/exp_A4_hard_instance.py
python experiments/exp_A4b_instances.py
python experiments/make_digest.py
cd theory && pdflatex main.tex && bibtex main && pdflatex main.tex && pdflatex main.tex
python report/build.py                  # make_restated.py, main.pdf, synthesis.pdf
\end{verbatim}
\end{small}
\nolinkurl{experiments/run_all.py} runs the scripts of all the experiments, those of hypotheses B3, B2, B1, A3, A2 and
A4 (\nolinkurl{exp_A4_hard_instance.py} and \nolinkurl{exp_A4b_instances.py}), of Section~\ref{sec:exp-constants} and of
Section~\ref{sec:exp-validation}, together with \nolinkurl{check_B2_fp_heuristic.py}, the sanity check of
\resref{rem:B2-gd-heuristic} (it writes \nolinkurl{results/B2_fp_check_summary.json}), and then
\nolinkurl{experiments/make_digest.py}; with \texttt{-{}-quick} it runs a smoke test, whose figures get a
\nolinkurl{_quick} suffix, and skips the digest. Results are cached in \nolinkurl{results/*.pkl} (delete a cache to
recompute its experiment) and summarised in \nolinkurl{results/*_summary.json}, which
\nolinkurl{experiments/make_digest.py} collects in \nolinkurl{results/digest.md}; every number quoted in this report is
read from these files or computed from them and from the cached runs. The statements restated in
Section~\ref{sec:theory} are not typed in the report: \nolinkurl{report/make_restated.py} copies them from the theory
notes into \nolinkurl{report/restated/} (one file per labelled result of the notes, and
\nolinkurl{report/restated/names.tex} with the type, appendix and name of each), so that the notes remain their only
source. It must be run before \nolinkurl{report/main.tex} and \nolinkurl{report/synthesis.tex}, the standalone extract
of Section~\ref{sec:theory}, are compiled. \nolinkurl{report/build.py} runs it, then compiles \nolinkurl{main.tex} and
\nolinkurl{synthesis.tex} in this order (pdflatex, bibtex and pdflatex twice for each), since the cross-references of
the extract are read from \nolinkurl{main.aux}.

\paragraph{Tests.} The test suite (\nolinkurl{tests/test_vrlab.py}) checks the gradients and Hessians of logistic
regression against finite differences and its per-sample smoothness constants; the Newton solver, on a small problem
($\norm{\Fp(\etastar)}<10^{-12}$) and on the problem $\kappa=10$, $\alpha=1$ of Section~\ref{sec:exp-B3}, where the
corrected line search must reach $\norm{\Fp(\etastar)}\le10^{-15}$ within ten iterations; the constants and the exact
minimizers of the instances (H), (S) and (A), including the geometric minimizers of (S) and (A); that the exact
conditional expectations used in the tightness analysis coincide with the average over the $n$ possible indices of the
realized next state, for SAGA in the Euclidean metric, in a dense and in a diagonal quadratic metric and on (S), and for
the Lyapunov function of L-Katyusha (both outcomes of the coin included); that the tightness ratios of I1--I5 lie in
$(0,1]$ at a random state; that the attribution chains of orders~A and~B (\resref*{prop:B3-attribution}) are monotone
and connect the exact contraction to the certified rate; that, for a pair $(\gamma,\Delta)$ other than the textbook
one, the a-posteriori factors of the trajectory-dependent bound (\resref*{prop:B3-traj}) dominate the exact one-step
contractions; that the exact one-step contraction of the Lyapunov function of L-Katyusha stays below the largest
coefficient of the one-step inequality of \resref{thm:lk}, at a random state and along a short run; that the mean of
$\norm{\theta_t-\etastar}^2$ over six SAGA runs on a small problem stays below the bound of Proposition~\ref{prop:59};
that SAGA, SAGA in a diagonal and in a dense quadratic metric, L-SVRG, L-Katyusha, GD and AGD converge on small
problems, and so do L-Katyusha with $\theta_2=0$ on two well-conditioned problems ($\kappa=5$ and $20$) and the naive
scheme (Nesterov momentum with the SVRG estimator) on an easy problem, whose gradient count is also checked; that the
error of SGD decreases; and the log--log slope fit with its interval. The divergence of L-Katyusha with $\theta_2=0$ on
ill-conditioned instances is observed in the experiments only.

\paragraph{Experiments.} Seeds: the logistic data sets are generated with seed~$0$ (for each $\alpha$, one data set
shared by all values of $\kappa$, which only set $\lambda$), the random block quadratics (R) with the sampling seed (one
instance per seed), and the chain instances (H), (S) and (A) are deterministic; the algorithms are run with seeds
$0, \dots, S-1$ (the number of seeds $S$ is stated with each figure). Every ``expectation'' of a one-step quantity is
computed exactly by enumerating the sampled index; curves labelled ``mean over seeds'' are empirical averages over the
$S$ runs. Trajectories are analysed only above the rounding floor ($10^{-18}$ on $V_t$ for SAGA and on
$\norm{\theta_t-\etastar}^2$ in Section~\ref{sec:exp-B2}, a floor set by the precision of $\etastar$; $10^{-10}$ for the
Lyapunov function of L-Katyusha, whose $W$-term amplifies the rounding error of $F(w) - F(\etastar)$ by the factor
$\theta_2(1+\theta_1)/(p\theta_1)$, from about $250$ to about $2400$ in our runs). Running times on one core range from
tens of seconds per logistic configuration of Section~\ref{sec:exp-B3} to several hours for each $\kappa$ sweep.

\paragraph{CPU-time budget of the $\kappa$ sweeps.} In the sweeps of Section~\ref{sec:exp-A4}, every solver run is capped
by \nolinkurl{VRLAB_MAX_TIME} seconds of CPU time, measured with \texttt{time.process\_time} (default $600$; a suspended
machine consumes no budget), and a run that reaches the cap does not enter the fits. The reported sweeps were run with
single-threaded BLAS and a budget of $4000$\,s per run, raised to $8000$\,s for the last runs of the A4b sweep, among them
SAGA at the nominal $\kappa=10^6$ on (S); every run reached its target, except the runs of L-Katyusha without negative
momentum, which diverge for $\kappa\ge10^3$. The grid search of the naive scheme gives each grid run $45$\,s of CPU time;
the grid runs have the same cost per iteration, so a run cut by this budget had already used more gradients than the
selected setting needed in total, and the cap does not affect the selection.

\paragraph{Minimizers.} For the quadratic instances, $\etastar$ is obtained by a direct solve. For logistic regression,
the cached results were computed with a damped Newton method whose Armijo test, decided by the rounding error of $F$ near
the optimum, stalled at $\norm{\Fp(\etastar)}=1.3\cdot10^{-12}$ on one problem ($\kappa=10$, $\alpha=1$; at most
$9\cdot10^{-16}$ on the fourteen other problems of Section~\ref{sec:exp-B3}). By strong convexity the squared error of
$\etastar$ is at most $\norm{\Fp(\etastar)}^2/\mu^2$: $2.1\cdot10^{-21}$ on that problem, where the corrected solver
moves $\etastar$ by a squared distance of $1.3\cdot10^{-21}$, and $7\cdot10^{-20}$ at most over the fifteen problems,
attained at $\kappa=10^5$, $\alpha=0$, whose runs stay above $10^{-13}$. On the problems whose runs reach the analysis
floor, these errors change $\norm{\theta_t-\etastar}^2=10^{-18}$ by at most about $10\%$. The line search has since been
corrected in \nolinkurl{vrlab/solvers.py} (near the optimum, a step is accepted as soon as it decreases $\norm{\Fp}$);
the results reported here were computed before this correction.

\paragraph{Figures.} Colours follow a fixed, colour-blind-safe categorical palette: a given algorithm or a given
inequality keeps its colour across all figures; sequential magnitudes use a single-hue ramp.
